\documentclass[10pt,a4paper]{amsart}
\usepackage[margin=19mm]{geometry}
\usepackage{amsmath,amssymb,mathtools}
\usepackage[scr=rsfs,cal=euler]{mathalfa}
\usepackage{xfrac}
\usepackage[unicode,hidelinks,bookmarks=false]{hyperref}
\newcommand{\ie}{i.e.\ }
\newcommand{\cf}{cf.\ }
\newcommand{\resp}{respectively\ }
\newcommand{\Subsection}{\subsection}
\providecommand{\nobreakdash}{\nobreak\hbox{-}}
\setlength{\emergencystretch}{2em}
\multlinegap0pt
\usepackage{bm}
\usepackage{bbm}
\usepackage{dsfont}
\usepackage{xspace}
\usepackage{cancel}
\usepackage{mathtools}
\usepackage{amsthm}
\makeatletter
\@ifundefined{theorem}{\newtheorem{theorem}{Theorem}}{}
\@ifundefined{lemma}{\newtheorem{lemma}[theorem]{Lemma}}{}
\@ifundefined{remark}{\newtheorem{remark}[theorem]{Remark}}{}
\makeatother
\newcommand{\C}{\mathbb C}
\newcommand{\Z}{\mathbb Z}
\newcommand{\coloneq}{\coloneqq}
\newcommand{\de}{\mathrm{d}}
\newcommand{\g}{\mathfrak{g}}
\newcommand{\im}{\operatorname{Im}}
\newcommand{\re}{\operatorname{Re}}
\title[Asymptotics of Riemannian Lie groups with nilpotency step 2]{Asymptotics of Riemannian Lie groups with nilpotency step 2}
\author{Enrico Le Donne, Luca Nalon, Sebastiano Nicolussi Golo, Seung-Yeon Ryoo}
\date{Source: arXiv:2503.00560v2 (CC BY 4.0)}
\begin{document}
\maketitle

\noindent\emph{Source and licence.} Asymptotics of Riemannian Lie groups with nilpotency step 2. Version arXiv:2503.00560v2 (CC BY 4.0), distributed under the Creative Commons Attribution 4.0 International licence, as recorded in the arXiv licence metadata for this exact version and shown on the abstract page https://arxiv.org/abs/2503.00560v2. Full rights notice: licenses/arxiv-2503.00560-CC-BY-4.0.txt.\par\smallskip

\noindent\textit{Editorial scope.} Counted unit: the Lemma whose source label is \texttt{lemma lemma} in the article (source0.tex lines 528--538, source label \texttt{lemma lemma}), delivered together with the article's own complete original proof of it (source0.tex lines 540--559, one \texttt{proof} environment of 20 lines). No mathematics is written, repaired or re-derived: statement and proof are the article's own text line for line. Required context retained uncounted: def G (lines 374--376); formula integral (lines 396--398); def Gc (lines 513--519); def\_f\_prop (lines 522--525); complex\_endpoint (lines 569--571). Every definition, display and label of those lines is retained unchanged. Omitted from this package and not redistributed: 1--373, 377--395, 399--512, 520--521, 526--527, 539--539, 560--568, 572--1989. Nothing inside a retained excerpt is omitted or altered: every retained body is delivered exactly as the article's lines are written, so no omission receipt of any kind accompanies this package. Citations: the retained excerpts carry no citation command at all, so no bibliography entry of the article is transcribed and no reference is redistributed. Reference adaptation: none was needed and none was made; every reference inside the delivered unit resolves inside it, so no unresolved reference or citation is produced. Printed slips kept literal: no printed word, symbol, hypothesis or number of the counted statement or of its proof was altered. Layout and class: the article's own class is \texttt{amsart} with packages including amssymb, amsmath, amsthm, amsfonts, mathtools, lmodern, mathrsfs, fontenc, inputenc, babel, geometry, graphicx, braket, csquotes; the wrapper is the lane's pinned \texttt{amsart} editorial wrapper at 10pt on A4, which restates the theorem-like environments the retained text needs and adds the article's own macro definitions that the retained text needs (\texttt{\textbackslash C}, \texttt{\textbackslash Z}, \texttt{\textbackslash coloneq}, \texttt{\textbackslash de}, \texttt{\textbackslash g}, \texttt{\textbackslash im}, \texttt{\textbackslash re}), with extra wrapper packages (bm, bbm, dsfont, xspace, cancel, mathtools). The delivered numbering is the unit's own. Assets: no retained span contains a figure environment, a graphics-inclusion command, a PDF-page inclusion, a table or a TikZ picture; no image file is copied into this package, the package declares no asset dependency and no image file is redistributed with it. Licence: version arXiv:2503.00560v2 of this article is distributed under the Creative Commons Attribution 4.0 International licence, as recorded in the arXiv licence metadata for this exact version and shown on the abstract page https://arxiv.org/abs/2503.00560v2. A full rights notice is delivered at licenses/arxiv-2503.00560-CC-BY-4.0.txt. Characters: every retained excerpt is pure ASCII, so the delivered engine, XeLaTeX with its default Latin Modern text font, reports no missing character.
\begin{equation} \label{def G}
 G:= (\g,*), \quad \text{where} \quad x*y \coloneq x+y+\tfrac{1}{2}[x,y], \qquad \text{ for all } x,y \in \g.
\end{equation}

	\begin{equation} \label{formula integral}
		\gamma_u(1) = \int^1_0 u(t) \, \de t + \frac{1}{2} \int^1_0 \left[\int^t_0 u(s) \, \de s, u(t) \right] \, \de t. 
	\end{equation}

\begin{remark}
 We point out that the complexification $\g_\C$ of a real Lie algebra $\g$ is also a real Lie algebra. Moreover, $\g$ is a Lie subalgebra of $\g_\C$. If $\g$ is 2-step nilpotent, then also $\g_\C$ is 2-step nilpotent. If a Lie group $G$ is as in \eqref{def G}, then we define 
 \begin{equation} \label{def Gc}
 G_\C:= (\g_\C,*), \quad \text{where} \quad x*y \coloneq x+y+\tfrac{1}{2}[x,y], \qquad \text{ for all } x,y \in \g_\C.
\end{equation}
Clearly, $G$ is a Lie subgroup of $G_\C$.
\end{remark}

\begin{gather} 
 f_n \colon [0,1] \to \C, \quad t \mapsto f_n(t):= e^{2\pi i n t}, \label{def_f}\\
 \text { for which }\int^1_0 f_n(t)f_m(t) \, \de t = \delta_{n,-m}, \quad \text{for every $n,m \in \Z$.} \label{def_f_prop}
\end{gather}

\begin{equation} \label{complex_endpoint}
 \gamma_v(1) = \sum_{n \in E} \frac{1}{4 \pi n} [\im(c_n),\re(c_n)] \quad \text{and} \quad \int^1_0 \rho(v(t),v(t)) \, \de t = \frac{1}{2} \sum_{n \in E} \rho(c_n,c_n).
\end{equation}

\begin{lemma} \label{lemma lemma}
 Let $\g$ be a 2-step nilpotent Lie algebra, and fix an inner product $\rho$ on a subspace $\Delta \subseteq \g$. Let $G_\C$ be as in \eqref{def Gc}, and consider $v \in L^2([0,1],\Delta_\C)$ such that
 \begin{equation} \label{def_xi_lemma}
 v(t) = \sum_{n \in E} c_nf_n(t), 
 \end{equation}
 for some finite subset of non-zero integers $E$, and some $\{c_n\}_{n\in E} \subseteq \Delta_\C$. Then
 \begin{equation} \label{complex_endpoint_1}
 \gamma_v(1) = \sum_{n \in E} \frac{1}{4 \pi i n} [c_n,c_{-n}] \quad \text{and} \quad \int^1_0 \rho(v(t),v(t)) \, \de t
 = \sum_{n \in E} \rho(c_n,c_n).
\end{equation}
\end{lemma}

\begin{proof}
 First, for every $t \in [0,1]$, we compute
 \begin{eqnarray} \label{int_xi_t}
 \int^t_0 v(s) \, \de s &\stackrel{\eqref{def_xi_lemma}}{=}& \int^t_0 \sum_{n \in E} c_nf_n(t) \, \de t = \sum_{n \in E} \frac{c_n}{2 \pi i n}(f_n(t)-1).
 \end{eqnarray}
 where in the second equality we used that $0 \notin E$. In particular, we obtain $\int^1_0 v(t) \, \de t = 0$. Thus, we compute
 \begin{eqnarray*}
 \gamma_v(1) &\stackrel{\eqref{formula integral}}{=}& \int^1_0 v(t) \, \de t + 
 \frac{1}{2} \int^1_0 \left[\int^t_0 v(s) \, \de s, v(t) \right] \, \de t \\ &\stackrel{\eqref{def_xi_lemma}\eqref{int_xi_t}}{=}& \frac{1}{2} \int^1_0 \left[ \sum_{n \in E} \frac{c_n}{2 \pi i n}(f_n(t)-1), \sum_{m \in E} c_mf_m(t)\right] \, \de t \\
 &=& \sum_{m,n \in E} \frac{1}{4\pi i n} [c_n,c_m] \int^1_0 (f_n(t)-f_0(t))f_m(t) \, \de t \\
 &\stackrel{\eqref{def_f_prop}}{=}& \sum_{n,m \in E} \frac{1}{4\pi in} [c_n,c_{m}]\delta_{n,-m} = \sum_{n \in E} \frac{1}{4\pi in} [c_n,c_{-n}]. 
 \end{eqnarray*} 
 We then verified the first identity in \eqref{complex_endpoint}. For the second one, we compute 
 \begin{eqnarray*}
 \int^1_0 \rho(v(t),v(t)) \, \de t &\stackrel{\eqref{def_xi_lemma}}{=}& \int^1_0 \rho\left(\sum_{n \in E} c_nf_n(t),\sum_{m \in E} c_mf_m(t)\right) \, \de t \\
 &=& \sum_{m,n \in E} \rho(c_n,c_m)\int^1_0 f_n(t)\overline{f_m}(t) \, \de t \\
 &\stackrel{\eqref{def_f_prop}}{=}& \sum_{m,n \in E} \rho(c_n,c_m) \delta_{n,m} \, \de t = \sum_{n \in E} \rho(c_n,c_n),
 \end{eqnarray*}
 where we used that $\rho$ is Hermitian and that $\overline{f_m}(t)=f_{-m}(t)$, for every $m \in \Z$ and $t \in [0,1]$.
 \end{proof}

\end{document}
